% SC4022 --- Chapter 2: Centrality Measures (0->100 ground-up)
\chapter{Centrality Measures}
\label{ch:centrality}
\sloppy

\begin{quote}
\emph{``Who matters most?''} --- In a friendship network, popularity differs from being a bridge; in the Web, a page linked by important pages is important. Different notions of ``central'' answer different questions.
\end{quote}

\noindent\fbox{\parbox{\dimexpr\linewidth-2\fboxsep-2\fboxrule}{%
\textbf{From zero: no prior centrality needed.} We assume only Ch.~1 (graph, degree, path, distance). Each measure is motivated by a story---counting friends, being near everyone, lying on many paths, being endorsed by the important---then defined, computed, and compared.}}
\vspace{0.4em}

\section*{Learning Outcomes}
After this chapter you will be able to:
\begin{enumerate}[label=(L\arabic*)]
    \item explain why degree alone misranks nodes when position matters;
    \item compute closeness and betweenness centralities on small graphs by hand;
    \item define eigenvector centrality and derive it from the adjacency spectrum;
    \item formulate PageRank as a random walk with teleportation and compute its stationary distribution;
    \item compare centrality rankings (correlation, rank disagreement) and select the right measure for a task.
\end{enumerate}

% ============================================================
\section{Why Not Just Degree?}
\label{sec:why-not-degree}
\paragraph{Intuition.} In a star, the centre has degree $5$ and every leaf has degree $1$---degree ranks them correctly. But attach two stars by a single bridge node (Fig.~\ref{fig:bridge-vs-hub}). The bridge has degree $2$, far below each star centre ($5$), yet removing the bridge disconnects the network while removing a leaf does nothing. ``Important'' depends on \emph{where} you sit, not just how many neighbours you have.

Consider also ``being close to everyone'': a node in the middle of a chain reaches all others in few hops, while a leaf at the end needs many hops, even if both have degree $2$ and $1$ respectively. And ``being endorsed by the important'': a junior researcher cited by two Nobel laureates may outrank one cited by ten undergraduates. These stories give closeness, betweenness, and eigenvector ideas.

\begin{figure}[h]\centering
\begin{tikzpicture}[scale=0.82,every node/.style={font=\scriptsize}]
% left star
\node[circle,draw,thick,fill=blue!18,minimum size=0.7cm] (c1) at (0,1.5) {H1};
\foreach \i/\a in {1/60,2/130,3/200,4/270,5/340} {\node[circle,draw,fill=blue!8,minimum size=0.52cm] (l1\i) at ($ (c1)+(\a:1.1) $) {L\i}; \draw[thick] (c1)--(l1\i);}
% bridge
\node[circle,draw,thick,fill=orange!18,minimum size=0.75cm] (br) at (3.0,1.5) {B};
\draw[thick,orange!70!black,line width=1.1pt] (c1)--(br);
% right star
\node[circle,draw,thick,fill=green!16,minimum size=0.7cm] (c2) at (6.0,1.5) {H2};
\foreach \i/\a in {1/60,2/130,3/200,4/270,5/340} {\node[circle,draw,fill=green!8,minimum size=0.52cm] (l2\i) at ($ (c2)+(\a:1.1) $) {R\i}; \draw[thick] (c2)--(l2\i);}
\draw[thick,orange!70!black,line width=1.1pt] (c2)--(br);
\node[font=\tiny,fill=orange!10,rounded corners=2pt,inner sep=2pt] at (3.0,0.05) {bridge: $k=2$ but cuts network};
\node[font=\tiny,fill=blue!10,rounded corners=2pt,inner sep=2pt] at (0,0.05) {hub $k=6$};
\node[font=\tiny,fill=green!10,rounded corners=2pt,inner sep=2pt] at (6.0,0.05) {hub $k=6$};
\node[font=\tiny,draw=gray!40,fill=yellow!8,rounded corners=2pt,inner sep=3pt] at (3.0,-0.65) {\textcolor{blue!60!black}{$\blacksquare$ left star} \quad \textcolor{orange!70!black}{$\blacksquare$ bridge} \quad \textcolor{green!50!black}{$\blacksquare$ right star}};
\end{tikzpicture}
\caption{Degree misranks when bridging matters. Both hubs have degree 6, the bridge has degree 2, yet the bridge lies on every path between left and right halves. Betweenness captures this; degree does not.}
\label{fig:bridge-vs-hub}
\end{figure}

\begin{definition}[Degree centrality]\label{def:deg-cent}
For node $v$, $C_D(v)=k_v$ (or normalised $k_v/(n-1)$). It is local---cost $O(m)$ for all nodes---and a useful baseline, but blind to global position.
\end{definition}

% ============================================================
\section{Distance-Based: Closeness and Betweenness}
\label{sec:closeness-betweenness}
\paragraph{Intuition for closeness.} If you must send a message to everyone, you want small total (or average) distance. A centre of a chain is closer on average than an endpoint.

\begin{definition}[Closeness and harmonic]\label{def:closeness}
Let $d(v,u)$ be shortest-path distance. Then
\begin{equation}
C_C(v)=\frac{n-1}{\sum_{u\neq v}d(v,u)}\qquad\text{(normalised closeness)}.
\end{equation}
If the graph is disconnected we use \emph{harmonic centrality} $C_H(v)=\sum_{u\neq v}1/d(v,u)$ with $1/\infty:=0$, which handles $\infty$ gracefully.
\end{definition}

\paragraph{Intuition for betweenness.} Stand on a node and count how many shortest paths you intercept. A bridge between communities intercepts many; a leaf intercepts none.

\begin{definition}[Betweenness]\label{def:betweenness}
Let $\sigma_{st}$ be the number of shortest $s\to t$ paths and $\sigma_{st}(v)$ those passing through $v$ ($s\neq v\neq t$). Then
\begin{equation}
C_B(v)=\sum_{s\neq v\neq t}\frac{\sigma_{st}(v)}{\sigma_{st}}.
\end{equation}
Normalise by $\binom{n-1}{2}$ for undirected graphs to get $[0,1]$. Brandes' algorithm computes all $C_B$ in $O(nm)$ (unweighted) via BFS from each source.
\end{definition}

\begin{example}[Hand computation on a 5-node chain]\label{ex:hand-betweenness}
Chain $1$--$2$--$3$--$4$--$5$. Distances from node 3: $1,1,2,2$ so $\sum d=6$ and $C_C(3)=4/6=0.667$. Endpoint 1 has distances $1,2,3,4$, sum $10$, $C_C(1)=0.4$---less central. Betweenness: node 3 lies on paths $1\to4,1\to5,2\to4,2\to5$ (4 pairs) plus symmetric, giving $C_B(3)=6$ unnormalised (highest); endpoints have $0$.
\end{example}

\begin{figure}[h]\centering
\begin{tikzpicture}[scale=0.82,every node/.style={font=\scriptsize}]
% chain positions
\foreach \i/\x in {1/0,2/2,3/4,4/6,5/8} {\node[circle,draw,thick,fill=blue!12,minimum size=0.7cm] (n\i) at (\x,1.2) {\i};}
\foreach \i/\j in {1/2,2/3,3/4,4/5} \draw[thick] (n\i)--(n\j);
% betweenness bars below
\foreach \i/\h/\c in {1/0.25/blue!12,2/0.65/orange!16,3/0.95/red!18,4/0.65/orange!16,5/0.25/blue!12} {
\draw[fill=\c,draw=gray!60] (n\i.south) ++(-0.3,-0.18) rectangle +(0.6,\h);
\node[font=\tiny] at ($(n\i.south)+(0,-0.05)$) {\h};}
\node[font=\tiny,fill=yellow!10,rounded corners=2pt,inner sep=2pt] at (4,2.35) {chain: centre is closest and most between};
\node[font=\tiny] at (4,-0.55) {bars = betweenness (taller = more paths through)};
\node[font=\tiny,draw=gray!40,fill=yellow!8,rounded corners=2pt,inner sep=3pt] at (4,-1.15) {\textcolor{red!70!black}{$\blacksquare$ highest $C_B$} \quad \textcolor{orange!60!black}{$\blacksquare$ medium} \quad \textcolor{blue!60!black}{$\blacksquare$ leaf $C_B=0$}};
% second row: star
\begin{scope}[yshift=-2.8cm]
\node[circle,draw,thick,fill=violet!16,minimum size=0.7cm] (s0) at (4,1.2) {0};
\foreach \i/\a in {1/30,2/90,3/150,4/210,5/270,6/330} {\node[circle,draw,fill=violet!8,minimum size=0.55cm] (s\i) at ($(s0)+(\a:1.2)$) {\i}; \draw[thick] (s0)--(s\i);}
\node[font=\tiny,fill=violet!10,rounded corners=2pt,inner sep=2pt] at (4,2.75) {star: centre $C_C=1$, $C_B=15$; leaves $0$};
\end{scope}
\end{tikzpicture}
\caption{Closeness and betweenness on a chain and a star. On the chain, the middle node maximises both; on the star, the centre dominates utterly. Bars sketch unnormalised $C_B$.}
\label{fig:close-between}
\end{figure}

% ============================================================
\section{Eigenvector and PageRank}
\label{sec:eigen-pagerank}
\paragraph{Intuition for eigenvector.} ``Important if linked by important nodes.'' This is self-referential: let $x_v$ be importance of $v$; then $x_v\propto\sum_{u\in N(v)}x_u$. In matrix form $Ax=\lambda x$---$x$ is an eigenvector of $A$.

\begin{definition}[Eigenvector centrality]\label{def:eigen}
Let $A$ be the adjacency matrix. The \emph{eigenvector centrality} is the entrywise non-negative eigenvector for the largest eigenvalue $\lambda_1$:
\begin{equation}
x_v=\frac1{\lambda_1}\sum_{u}A_{vu}x_u,\quad \|x\|_2=1,\;x_v\ge0.
\end{equation}
Perron--Frobenius guarantees existence for connected (irreducible) graphs. Power iteration $x^{(t+1)}\propto Ax^{(t)}$ converges to it.
\end{definition}

\paragraph{Intuition for PageRank.} Eigenvector fails on directed graphs with sinks (a node with out-degree 0 traps importance). PageRank imagines a random surfer: with probability $\alpha$ follow a random out-link, with probability $1-\alpha$ teleport to a random node. This makes the walk irreducible and aperiodic, so a unique stationary distribution exists---that distribution \emph{is} PageRank.

\begin{definition}[PageRank]\label{def:pagerank}
For directed graph with out-degree $k^{\text{out}}_u$, transition $P_{uv}=A_{uv}/k^{\text{out}}_u$ (or $1/n$ if $k^{\text{out}}_u=0$). With damping $\alpha\in(0,1)$ (usually $0.85$) and teleportation matrix $J_{uv}=1/n$,
\begin{equation}
G=\alpha P+(1-\alpha)J,\qquad \pi=\pi G,\quad \sum_v\pi_v=1,\;\pi_v\ge0.
\end{equation}
$\pi$ is the PageRank vector; $\pi_v$ is the long-run fraction of time at $v$.
\end{definition}

\begin{figure}[h]\centering
\begin{tikzpicture}[scale=0.82,every node/.style={font=\scriptsize}]
\node[circle,draw,thick,fill=green!15,minimum size=0.8cm] (a) at (0,1.3) {A};
\node[circle,draw,thick,fill=green!15,minimum size=0.8cm] (b) at (2.4,1.8) {B};
\node[circle,draw,thick,fill=green!15,minimum size=0.8cm] (c) at (2.4,0.4) {C};
\node[circle,draw,thick,fill=green!15,minimum size=0.8cm] (d) at (4.8,1.3) {D};
\draw[->,thick,>=Stealth,blue!60!black] (a) to (b); \draw[->,thick,>=Stealth,blue!60!black] (b) to (d);
\draw[->,thick,>=Stealth,blue!60!black] (a) to (c); \draw[->,thick,>=Stealth,blue!60!black] (c) to (d);
\draw[->,thick,>=Stealth,blue!60!black] (d) to[bend left=22] (a);
\draw[->,thick,dashed,orange!70!black] (b) to[bend left=18] (a);
\draw[->,thick,dashed,orange!70!black] (c) to[bend right=18] (a);
\node[font=\tiny,fill=blue!8,rounded corners=2pt,inner sep=2pt] at (2.4,2.45) {solid = follow link ($\alpha$)};
\node[font=\tiny,fill=orange!10,rounded corners=2pt,inner sep=2pt] at (2.4,-0.45) {dashed = teleport ($1-\alpha$)};
% bar convergence on right
\begin{scope}[xshift=7.0cm]
\foreach \i/\h in {1/0.45,2/0.62,3/0.71,4/0.76,5/0.78} {
\draw[fill=violet!\the\numexpr 8+6*\i\relax,draw=gray!60] (0.5*\i,0) rectangle +(0.35,\h);}
\node[font=\tiny] at (1.5,-0.3) {power iteration $\to\pi$};
\node[font=\tiny,rotate=90] at (-0.25,0.5) {$\pi_A$};
\end{scope}
\node[font=\tiny,draw=gray!40,fill=yellow!8,rounded corners=2pt,inner sep=3pt] at (2.4,-1.05) {\textcolor{blue!60!black}{$\blacksquare$ link} \quad \textcolor{orange!70!black}{$\blacksquare$ teleport} \quad \textcolor{violet!60!black}{$\blacksquare$ $G=\alpha P+(1-\alpha)J$}};
\end{tikzpicture}
\caption{PageRank as a random walk with teleportation. Solid arrows are link-following; dashed are random jumps that guarantee ergodicity. Right: power iteration converges to the stationary distribution $\pi$.}
\label{fig:pagerank}
\end{figure}

\begin{lstlisting}[language=Python, caption={Centrality zoo on Zachary karate club: degree, closeness, betweenness, eigenvector, PageRank.}, label={lst:centrality-zoo}]
import networkx as nx

G = nx.karate_club_graph()
# All centralities (normalised where applicable)
deg = nx.degree_centrality(G)
clo = nx.closeness_centrality(G)
bet = nx.betweenness_centrality(G, normalized=True)
eig = nx.eigenvector_centrality(G, max_iter=500)
pr  = nx.pagerank(G, alpha=0.85)

def top(d, k=3): return sorted(d, key=d.get, reverse=True)[:k]
print("degree:     ", top(deg), [round(deg[v],2) for v in top(deg)])
print("closeness:  ", top(clo), [round(clo[v],2) for v in top(clo)])
print("betweenness:", top(bet), [round(bet[v],3) for v in top(bet)])
print("eigenvector:", top(eig), [round(eig[v],3) for v in top(eig)])
print("pagerank:   ", top(pr),  [round(pr[v],3)  for v in top(pr)])
# Spearman correlation: do rankings agree?
import scipy.stats as st
print("deg vs bet rho:", round(st.spearmanr(list(deg.values()),
      list(bet.values())).correlation, 2))
\end{lstlisting}

\begin{lstlisting}[language=Python, caption={Betweenness reveals bridges: star-of-stars where degree misleads.}, label={lst:bridge-demo}]
import networkx as nx
# Two cliques joined by a bridge node
G = nx.Graph()
G.add_edges_from([(0,1),(0,2),(1,2)])       # left triangle
G.add_edges_from([(3,4),(3,5),(4,5)])       # right triangle
G.add_edge(2, 3)                             # bridge edge 2-3

print("degrees:", dict(G.degree()))
print("betweenness:", {v: round(b,3) for v,b in
      nx.betweenness_centrality(G).items()})
print("closeness:", {v: round(c,3) for v,c in
      nx.closeness_centrality(G).items()})
# Node 2 and 3 have modest degree but highest betweenness
\end{lstlisting}

% ============================================================
\section{Choosing a Centrality}
\label{sec:choosing}
No single measure is ``correct.'' Degree is cheap and local---good for immediate risk (hub vaccination). Closeness suits broadcast reach (place a warehouse). Betweenness suits bottlenecks (protect a bridge router). Eigenvector/PageRank suits endorsement networks (Web search, citation). A practical workflow: compute several, check rank correlation (Spearman $\rho$), inspect disagreements, and validate against an external outcome (e.g.\ does high PageRank predict traffic?).

% ============================================================
\section{Chapter Notes}
Centrality thinking began with Bavelas (1948) and Freeman (1977) formalising betweenness and closeness for social networks. Eigenvector centrality was refined by Bonacich (1972); PageRank (Brin \& Page, 1998) made it the backbone of Web search. Brandes (2001) gave the $O(nm)$ betweenness algorithm still in use. Modern caution: centralities are highly correlated on many real graphs---reporting all five without noting redundancy is a common pitfall.

% ============================================================
\section{Exercises}
\label{sec:ch02-exercises}
\begin{enumerate}[label=E2.\arabic*, leftmargin=*, itemsep=0.8em]
    \item \textbf{Degree vs position.} Draw the ``lollipop'': a $K_4$ (nodes $1$--$4$) plus a chain $4$--$5$--$6$. Rank all nodes by degree, then by betweenness (reason, no code). Which node is second-highest in betweenness despite modest degree? Explain.
    \item \textbf{Closeness by hand.} Compute $C_C(v)$ for every $v$ in the 4-node star (centre $1$ + leaves $2,3,4$). Then compute harmonic centrality and verify that ordering agrees but values differ.
    \item \textbf{Eigenvector iteration.} For the triangle graph $K_3$, start $x^{(0)}=[1,0,0]$ and do two steps of $x^{(t+1)}\propto Ax^{(t)}$ (normalise to sum 1). What limit do you expect and why? What happens if the graph has two components?
    \item \textbf{PageRank damping.} On a directed 3-cycle $1\to2\to3\to1$, what is PageRank for any $\alpha$? Add a dangling node $4$ with edge $3\to4$ (no out-links from $4$): how does $\pi_4$ vary with $\alpha$? Explain teleportation's role.
    \item \textbf{Correlated rankings.} Using NetworkX on the karate club, compute Spearman $\rho$ for every pair among degree, betweenness, and PageRank. Report the most and least correlated pair and inspect which nodes cause rank disagreements.
\end{enumerate}
% End of Chapter 2
